\capitulo{Variables conjuntas discretas: la tabla}{Leer la tabla, calcular probabilidades de eventos, sacar marginales, condicionar y decidir independencia.}{cap:discretas}
\meta{95}{Ej. 3.3 y 3.7 \textbullet{} Serie 1, 8 y 10}

\subsection{Qu\'e es una funci\'on de probabilidad conjunta}

En el Tema 2 un experimento contestaba \textbf{una} pregunta ($X$). Ahora contesta \textbf{dos} a la vez: $\cx{X}$ y $\cy{Y}$.
Por ejemplo, en una f\'abrica con dos l\'ineas: $\cx{X}$ = art\'iculos que produce la l\'inea uno en un d\'ia,
$\cy{Y}$ = art\'iculos que produce la l\'inea dos ese mismo d\'ia.

La \textbf{funci\'on de probabilidad conjunta} dice qu\'e tan probable es cada \emph{pareja} de resultados:
\[\fxy(\cx{x},\cy{y})=P(\cx{X=x},\ \cy{Y=y})\qquad\text{la coma se lee \guillemotleft y al mismo tiempo\guillemotright}\]

Cuando las variables son discretas, la conjunta es una \textbf{tabla}. \'Esta es la de la f\'abrica (Ej.~3.3 de tus apuntes)
y la vamos a usar en todo el cap\'itulo:

\begin{center}
\begin{tabular}{|C|C|C|C|}\hline
\hd{$y\,\backslash\,x$} & \hx{0} & \hx{1} & \hx{2}\\\hline
\hy{0} & 0.10 & \ce{0.20} & 0.20\\\hline
\hy{1} & 0.04 & 0.08 & 0.08\\\hline
\hy{2} & 0.06 & 0.12 & 0.12\\\hline
\end{tabular}
\end{center}

\textbf{C\'omo se lee una celda.} La celda marcada est\'a en la columna $\cx{x=1}$ y el rengl\'on $\cy{y=0}$:
$\fxy(1,0)=P(X=1,\,Y=0)=0.20$. \enpalabras{en el 20\,\% de los d\'ias la l\'inea uno produce 1 art\'iculo y, ese mismo d\'ia,
la l\'inea dos produce 0.}

\begin{idea}{Las dos condiciones de toda tabla conjunta}
\begin{enumerate}[label=\protect\circnum{\arabic*},leftmargin=10mm,labelsep=3mm]
\item Cada celda est\'a entre 0 y 1: \quad $0\le\fxy(x,y)\le1$.
\item Todas las celdas juntas suman 1: \quad $\displaystyle\sum_x\sum_y\fxy(x,y)=1$.
\end{enumerate}
Si en una tabla falta un valor (un \guillemotleft?\guillemotright), se encuentra con la condici\'on 2.
\end{idea}

\begin{ojo}
\mal{Suponer que los renglones siempre son $x$.}
\bien{Revisa el encabezado. En esta tabla los renglones son $\cy{y}$ y las columnas $\cx{x}$; en la Serie~3, problema~1,
es al rev\'es. Escribe arriba de tu hoja: \guillemotleft renglones = \ldots, columnas = \ldots\guillemotright.}
\end{ojo}

\subsection{Probabilidad de un evento: marcar y sumar}

Un evento es una condici\'on sobre las dos variables: $X>Y$, $X=Y$, $X+Y\le2$.
Su probabilidad es la suma de las celdas que cumplen la condici\'on.

\begin{resuelto}{Probabilidad de que la l\'inea uno produzca m\'as}{Apuntes \textbullet{} Ej.~3.3 a}
\begin{enunciado}
Con la tabla de la f\'abrica, calcular $P(X>Y)$.
\end{enunciado}
\begin{pasos}
\paso{Reviso celda por celda si se cumple $x>y$}
Voy por renglones. En $y=0$: ?`$0>0$? no; ?`$1>0$? s\'i; ?`$2>0$? s\'i.
En $y=1$: ?`$0>1$? no; ?`$1>1$? no; ?`$2>1$? s\'i. En $y=2$: ninguna cumple.
\paso{Marco las que cumplen}
\begin{center}
\begin{tabular}{|C|C|C|C|}\hline
\hd{$y\,\backslash\,x$} & \hx{0} & \hx{1} & \hx{2}\\\hline
\hy{0} & 0.10 & \ce{0.20} & \ce{0.20}\\\hline
\hy{1} & 0.04 & 0.08 & \ce{0.08}\\\hline
\hy{2} & 0.06 & 0.12 & 0.12\\\hline
\end{tabular}
\end{center}
\paso{Sumo las marcadas}
\[P(X>Y)=0.20+0.20+0.08=\res{0.48}\]
\end{pasos}
\porque{Cada celda es la probabilidad de una pareja distinta. Como no pueden ocurrir dos parejas a la vez, sus probabilidades se suman.}
\end{resuelto}

\begin{turno}
Con la misma tabla calcula \textbf{a)} $P(X=Y)$ \quad \textbf{b)} $P(X+Y\le2)$.
\comprueba{\textbf{a)} la diagonal: $0.10+0.08+0.12=0.30$.\quad
\textbf{b)} celdas $(0,0),(1,0),(2,0),(0,1),(1,1),(0,2)$: $0.10+0.20+0.20+0.04+0.08+0.06=0.68$.}
\end{turno}

\subsection{Marginales: quedarse con una sola variable}

A veces s\'olo importa una variable. ?`Cu\'al es la probabilidad de que la l\'inea uno produzca 1 art\'iculo,
\textbf{sin importar} qu\'e haga la l\'inea dos? Hay tres maneras de que $X=1$ ocurra (con $Y=0$, con $Y=1$ o con $Y=2$);
se suman las tres.

\begin{idea}{Marginal = sumar la variable que no me interesa}
\[f_X(\cx{x})=\sum_{\cy{y}}\fxy(\cx{x},\cy{y})\qquad\qquad f_Y(\cy{y})=\sum_{\cx{x}}\fxy(\cx{x},\cy{y})\]
En la tabla: la marginal de $\cx{X}$ es el \textbf{total de cada columna de $x$}; la marginal de $\cy{Y}$ es el
\textbf{total de cada rengl\'on de $y$}. Se llaman marginales porque se escriben al margen de la tabla.
\end{idea}

\begin{center}
\begin{tikzpicture}[x=1.7cm,y=-0.85cm,
  celda/.style={draw=tinta!30,minimum width=1.5cm,minimum height=0.72cm,font=\small},
  hX/.style={fill=colX,text=white,minimum width=1.5cm,minimum height=0.72cm,font=\sffamily\bfseries\small},
  hY/.style={fill=colY,text=white,minimum width=1.5cm,minimum height=0.72cm,font=\sffamily\bfseries\small}]
\foreach \v [count=\i from 0] in {0,1,2} \node[hX] at (\i,-1) {$x=\v$};
\foreach \v [count=\j from 0] in {0,1,2} \node[hY] at (-1,\j) {$y=\v$};
\foreach \row [count=\j from 0] in {{0.10,0.20,0.20},{0.04,0.08,0.08},{0.06,0.12,0.12}}
  \foreach \v [count=\i from 0] in \row \node[celda] at (\i,\j) {\v};
\foreach \v [count=\j from 0] in {0.50,0.20,0.30}{
  \draw[-{Stealth},colY,line width=1.1pt] (2.52,\j)--(3.2,\j);
  \node[celda,fill=colY!18,draw=colY,font=\bfseries\small] at (3.7,\j) {\v};}
\foreach \v [count=\i from 0] in {0.20,0.40,0.40}{
  \draw[-{Stealth},colX,line width=1.1pt] (\i,2.5)--(\i,3.15);
  \node[celda,fill=colX!16,draw=colX,font=\bfseries\small] at (\i,3.65) {\v};}
\node[font=\sffamily\bfseries\small,text=colY] at (3.7,-1) {$f_Y(y)$};
\node[font=\sffamily\bfseries\small,text=colX] at (-1,3.65) {$f_X(x)$};
\node[font=\small,text=colY,anchor=west] at (4.3,1) {sumo cada rengl\'on};
\node[font=\small,text=colX,anchor=north] at (1,4.15) {sumo cada columna};
\end{tikzpicture}
\end{center}

\begin{resuelto}{Marginales de la f\'abrica}{Apuntes \textbullet{} Ej.~3.3 b}
\begin{pasos}
\paso{Marginal de $X$: sumo cada columna}
\begin{align*}
f_X(0)&=0.10+0.04+0.06=0.20\\
f_X(1)&=0.20+0.08+0.12=0.40\\
f_X(2)&=0.20+0.08+0.12=0.40
\end{align*}
\paso{Marginal de $Y$: sumo cada rengl\'on}
\begin{align*}
f_Y(0)&=0.10+0.20+0.20=0.50\\
f_Y(1)&=0.04+0.08+0.08=0.20\\
f_Y(2)&=0.06+0.12+0.12=0.30
\end{align*}
\paso{Las escribo al margen de la tabla}
\begin{center}
\begin{tabular}{|C|C|C|C|C|}\hline
\hd{$y\,\backslash\,x$} & \hx{0} & \hx{1} & \hx{2} & \hy{$f_Y(y)$}\\\hline
\hy{0} & 0.10 & 0.20 & 0.20 & \my{0.50}\\\hline
\hy{1} & 0.04 & 0.08 & 0.08 & \my{0.20}\\\hline
\hy{2} & 0.06 & 0.12 & 0.12 & \my{0.30}\\\hline
\hx{$f_X(x)$} & \mx{0.20} & \mx{0.40} & \mx{0.40} & \textbf{1}\\\hline
\end{tabular}
\end{center}
\end{pasos}
\verifica{cada marginal suma 1 ($0.20+0.40+0.40=1$ y $0.50+0.20+0.30=1$). Si no, hay un error de suma.}
\enpalabras{$f_X(1)=0.40$ significa que la l\'inea uno produce 1 art\'iculo el 40\,\% de los d\'ias, haga lo que haga la l\'inea dos.}
\end{resuelto}

\subsection{Condicional: cuando ya s\'e cu\'anto vali\'o la otra}

Ahora me informan algo: \guillemotleft hoy la l\'inea dos produjo 1 art\'iculo\guillemotright{} ($\cy{Y=1}$). Con ese dato,
ya no me interesa toda la tabla: \textbf{s\'olo el rengl\'on $y=1$}. Pero ese rengl\'on suma $0.20$, no 1, as\'i que lo
\textbf{divido entre su total} para que vuelva a ser una distribuci\'on.

\begin{center}
\begin{tikzpicture}[x=1.7cm,y=-0.85cm,
  celda/.style={draw=tinta!30,minimum width=1.5cm,minimum height=0.72cm,font=\small}]
\node[font=\small,anchor=east] at (-0.6,0) {rengl\'on $\cy{y=1}$:};
\foreach \v [count=\i from 0] in {0.04,0.08,0.08} \node[celda,fill=marca] at (\i,0) {\v};
\node[font=\small,anchor=west] at (2.6,0) {suma $0.20=f_Y(1)$};
\draw[-{Stealth},tinta,line width=1.1pt] (1,0.55)--(1,1.45) node[midway,right,font=\small] {divido cada celda entre $0.20$};
\node[font=\small,anchor=east] at (-0.6,2) {$f_{X\mid Y=1}(x)$:};
\foreach \v [count=\i from 0] in {0.2,0.4,0.4} \node[celda,fill=verde!10,draw=verde] at (\i,2) {\v};
\node[font=\small,anchor=west] at (2.6,2) {ahora suma 1};
\end{tikzpicture}
\end{center}

\begin{idea}{Condicional = conjunta entre la marginal de lo que ya conozco}
\[
f_{X\mid y_0}(x\mid Y=y_0)=\frac{\fxy(x,y_0)}{f_Y(y_0)}
\qquad\qquad
f_{Y\mid x_0}(y\mid X=x_0)=\frac{\fxy(x_0,y)}{f_X(x_0)}
\]
Abajo va siempre la marginal de la variable \textbf{que est\'a despu\'es de la barra} (la que ya se conoce).
\end{idea}

\begin{resuelto}{Condicional de $X$ dado $Y=1$}{Apuntes \textbullet{} Ej.~3.3 c, d}
\begin{enunciado}
Determinar la funci\'on de probabilidad condicional de $X$ dado $Y=1$ y calcular $P(X<2\mid Y=1)$.
\end{enunciado}
\begin{pasos}
\paso{Identifico qu\'e se conoce}
Se conoce $Y=1$. Uso el rengl\'on $y=1$ y su total, $f_Y(1)=0.20$.
\paso{Divido cada celda del rengl\'on entre el total}
\begin{align*}
P(X=0\mid Y=1)&=\frac{\fxy(0,1)}{f_Y(1)}=\frac{0.04}{0.20}=0.2\\[2pt]
P(X=1\mid Y=1)&=\frac{\fxy(1,1)}{f_Y(1)}=\frac{0.08}{0.20}=0.4\\[2pt]
P(X=2\mid Y=1)&=\frac{\fxy(2,1)}{f_Y(1)}=\frac{0.08}{0.20}=0.4
\end{align*}
\paso{Presento la condicional como tabla}
\begin{center}
\begin{tabular}{|D|C|C|C|}\hline
\hx{$x$} & 0 & 1 & 2\\\hline
\hx{$f_{X\mid Y=1}$} & 0.2 & 0.4 & 0.4\\\hline
\end{tabular}
\end{center}
\verifica{$0.2+0.4+0.4=1$.}
\paso{Calculo la probabilidad pedida con esa tabla}
$X<2$ significa $X=0$ o $X=1$:
\[P(X<2\mid Y=1)=0.2+0.4=\res{0.6}\]
\end{pasos}
\end{resuelto}

\subsection{Independencia: ?`una variable informa sobre la otra?}

Compara lo que acabamos de obtener: la condicional $f_{X\mid Y=1}=(0.2,\,0.4,\,0.4)$ es \textbf{id\'entica} a la marginal
$f_X=(0.20,\,0.40,\,0.40)$. Saber que $Y=1$ no cambi\'o nada lo que se espera de $X$. Eso es independencia.

\begin{idea}{Independencia}
$X$ y $Y$ son independientes si y s\'olo si, \textbf{en todas las celdas},
\[\fxy(x,y)=f_X(x)\cdot f_Y(y)\]
\begin{itemize}
\item Para decir que \textbf{no} son independientes basta \textbf{una} celda donde falle.
\item Para decir que \textbf{s\'i} lo son hay que revisar \textbf{todas}.
\end{itemize}
\end{idea}

\begin{resuelto}{?`Son independientes las dos l\'ineas?}{Apuntes \textbullet{} Ej.~3.3}
\begin{pasos}
\paso{Tengo la tabla con sus m\'argenes}
$f_X=(0.20,\,0.40,\,0.40)$ y $f_Y=(0.50,\,0.20,\,0.30)$.
\paso{Construyo la tabla de productos \textnormal{(marginal de la columna por marginal del rengl\'on)}}
\begin{center}\small
\begin{tabular}{|C|D|D|D|}\hline
\hd{$f_X\cdot f_Y$} & \hx{$x=0$} & \hx{$x=1$} & \hx{$x=2$}\\\hline
\hy{$y=0$} & $0.2(0.5)=0.10$ & $0.4(0.5)=0.20$ & $0.4(0.5)=0.20$\\\hline
\hy{$y=1$} & $0.2(0.2)=0.04$ & $0.4(0.2)=0.08$ & $0.4(0.2)=0.08$\\\hline
\hy{$y=2$} & $0.2(0.3)=0.06$ & $0.4(0.3)=0.12$ & $0.4(0.3)=0.12$\\\hline
\end{tabular}
\end{center}
\paso{Comparo con la tabla original}
Las nueve celdas coinciden con la conjunta. \res{\text{S\'i son independientes}}
\end{pasos}
\end{resuelto}

\begin{receta}{cualquier problema con tabla}
\begin{pasosreceta}
\item \textbf{Ubica qui\'en es qui\'en:} qu\'e variable est\'a en renglones y cu\'al en columnas.
\item \textbf{Agrega los m\'argenes:} total de cada rengl\'on y de cada columna. Comprueba que todo sume 1.
\item \textbf{Evento} ($X>Y$, $X+Y=4$\ldots): marca las celdas que cumplen y s\'umalas.
\item \textbf{Condicional:} toma el rengl\'on o la columna del valor conocido y divide entre su total.
\item \textbf{Independencia:} compara cada celda con (su margen de rengl\'on)$\times$(su margen de columna).
\end{pasosreceta}
\end{receta}

\begin{resuelto}{Supermercado}{Serie 3 \textbullet{} problema 1}
\begin{enunciado}
$X_1$ = clientes en la fila de la caja r\'apida, $X_2$ = clientes en la fila de la caja superveloz.
\begin{center}\small
\begin{tabular}{|C|C|C|C|C|}\hline
\hd{$x_1\backslash\,x_2$} & \hy{0} & \hy{1} & \hy{2} & \hy{3}\\\hline
\hx{0} & 0.08 & 0.07 & 0.04 & 0.00\\\hline
\hx{1} & 0.06 & 0.15 & 0.05 & 0.04\\\hline
\hx{2} & 0.05 & 0.04 & 0.10 & 0.06\\\hline
\hx{3} & 0.00 & 0.03 & 0.04 & 0.07\\\hline
\hx{4} & 0.00 & 0.01 & 0.05 & 0.06\\\hline
\end{tabular}
\end{center}
a) $P(X_1=1,X_2=1)$.\quad b) $P(X_1=X_2)$.\quad c) $A$ = hay por lo menos dos clientes m\'as en una fila que en la otra; expresar $A$ y calcular $P(A)$.\quad
d) Probabilidad de que el total de clientes sea exactamente cuatro, y de que sea por lo menos cuatro.
\end{enunciado}
Aqu\'i los renglones son $\cx{x_1}$ y las columnas $\cy{x_2}$.

\inciso{a}{Exactamente un cliente en cada fila}
Es una sola celda: rengl\'on $x_1=1$, columna $x_2=1$. \quad $P(X_1=1,X_2=1)=\res{0.15}$

\inciso{b}{Mismo n\'umero de clientes en las dos filas}
\begin{pasos}
\paso{Marco las celdas con $x_1=x_2$ \textnormal{(la diagonal)}}
\begin{center}\small
\begin{tabular}{|C|C|C|C|C|}\hline
\hd{$x_1\backslash\,x_2$} & \hy{0} & \hy{1} & \hy{2} & \hy{3}\\\hline
\hx{0} & \ce{0.08} & 0.07 & 0.04 & 0.00\\\hline
\hx{1} & 0.06 & \ce{0.15} & 0.05 & 0.04\\\hline
\hx{2} & 0.05 & 0.04 & \ce{0.10} & 0.06\\\hline
\hx{3} & 0.00 & 0.03 & 0.04 & \ce{0.07}\\\hline
\hx{4} & 0.00 & 0.01 & 0.05 & 0.06\\\hline
\end{tabular}
\end{center}
\paso{Sumo}
\[P(X_1=X_2)=0.08+0.15+0.10+0.07=\res{0.40}\]
\end{pasos}

\inciso{c}{Por lo menos dos clientes m\'as en una fila que en la otra}
\begin{pasos}
\paso{Traduzco la frase a s\'imbolos}
\guillemotleft Dos m\'as en una que en la otra\guillemotright{} habla de la \textbf{diferencia} sin importar cu\'al fila es mayor,
y \guillemotleft por lo menos\guillemotright{} es $\ge$: \quad $A=\{\,|X_1-X_2|\ge2\,\}$.
\paso{Marco las celdas donde la diferencia es 2 o m\'as}
\begin{center}\small
\begin{tabular}{|C|C|C|C|C|}\hline
\hd{$x_1\backslash\,x_2$} & \hy{0} & \hy{1} & \hy{2} & \hy{3}\\\hline
\hx{0} & 0.08 & 0.07 & \ce{0.04} & \ce{0.00}\\\hline
\hx{1} & 0.06 & 0.15 & 0.05 & \ce{0.04}\\\hline
\hx{2} & \ce{0.05} & 0.04 & 0.10 & 0.06\\\hline
\hx{3} & \ce{0.00} & \ce{0.03} & 0.04 & 0.07\\\hline
\hx{4} & \ce{0.00} & \ce{0.01} & \ce{0.05} & 0.06\\\hline
\end{tabular}
\end{center}
\paso{Sumo}
\[P(A)=0.04+0.00+0.04+0.05+0.00+0.03+0.00+0.01+0.05=\res{0.22}\]
\end{pasos}

\inciso{d}{Total exactamente cuatro, y por lo menos cuatro}
\begin{pasos}
\paso{Exactamente cuatro: celdas con $x_1+x_2=4$}
Son $(1,3)$, $(2,2)$, $(3,1)$ y $(4,0)$:
\[P(X_1+X_2=4)=0.04+0.10+0.03+0.00=\res{0.17}\]
\paso{Por lo menos cuatro: agrego las celdas que suman 5, 6 y 7}
\begin{center}\small
\begin{tabular}{|C|C|C|C|C|}\hline
\hd{$x_1\backslash\,x_2$} & \hy{0} & \hy{1} & \hy{2} & \hy{3}\\\hline
\hx{0} & 0.08 & 0.07 & 0.04 & 0.00\\\hline
\hx{1} & 0.06 & 0.15 & 0.05 & \ce{0.04}\\\hline
\hx{2} & 0.05 & 0.04 & \ce{0.10} & \ce{0.06}\\\hline
\hx{3} & 0.00 & \ce{0.03} & \ce{0.04} & \ce{0.07}\\\hline
\hx{4} & \ce{0.00} & \ce{0.01} & \ce{0.05} & \ce{0.06}\\\hline
\end{tabular}
\end{center}
\[P(X_1+X_2\ge4)=\underbrace{0.17}_{\text{suman }4}+\underbrace{0.06+0.04+0.01}_{\text{suman }5}+\underbrace{0.07+0.05}_{\text{suman }6}+\underbrace{0.06}_{\text{suman }7}=\res{0.46}\]
\end{pasos}
\end{resuelto}

\begin{resuelto}{Restaurante}{Serie 3 \textbullet{} problema 8, incisos a, b, c}
\begin{enunciado}
$X$ = costo de la comida del hombre, $Y$ = costo de la comida de la mujer (7, 9 o 10 pesos).
\begin{center}\small
\begin{tabular}{|C|C|C|C|}\hline
\hd{$x\,\backslash\,y$} & \hy{7} & \hy{9} & \hy{10}\\\hline
\hx{7} & 0.05 & 0.05 & 0.10\\\hline
\hx{9} & 0.05 & 0.10 & 0.35\\\hline
\hx{10} & 0 & 0.20 & 0.10\\\hline
\end{tabular}
\end{center}
a) Marginales.\quad b) Probabilidad de que ambas comidas cuesten a lo sumo \$9.\quad c) ?`Son independientes?
\end{enunciado}

\inciso{a}{Marginales}
\begin{pasos}
\paso{Marginal de $X$ \textnormal{(aqu\'i $x$ est\'a en renglones)}: sumo cada rengl\'on}
$f_X(7)=0.05+0.05+0.10=0.20$;\quad $f_X(9)=0.05+0.10+0.35=0.50$;\quad $f_X(10)=0+0.20+0.10=0.30$.
\paso{Marginal de $Y$: sumo cada columna}
$f_Y(7)=0.05+0.05+0=0.10$;\quad $f_Y(9)=0.05+0.10+0.20=0.35$;\quad $f_Y(10)=0.10+0.35+0.10=0.55$.
\paso{Resultado}
\begin{center}\small
\begin{tabular}{|C|C|C|C|C|}\hline
\hd{$x\,\backslash\,y$} & \hy{7} & \hy{9} & \hy{10} & \hx{$f_X(x)$}\\\hline
\hx{7} & 0.05 & 0.05 & 0.10 & \mx{0.20}\\\hline
\hx{9} & 0.05 & 0.10 & 0.35 & \mx{0.50}\\\hline
\hx{10} & 0 & 0.20 & 0.10 & \mx{0.30}\\\hline
\hy{$f_Y(y)$} & \my{0.10} & \my{0.35} & \my{0.55} & \textbf{1}\\\hline
\end{tabular}
\end{center}
\end{pasos}

\inciso{b}{Ambas cuestan a lo sumo \$9}
\guillemotleft A lo sumo 9\guillemotright{} es $\le9$, para las dos: $X\le9$ \textbf{y} $Y\le9$. Son las cuatro celdas de la esquina superior izquierda:
\[P(X\le9,\,Y\le9)=0.05+0.05+0.05+0.10=\res{0.25}\]

\inciso{c}{?`Independientes?}
\begin{pasos}
\paso{Busco una celda sospechosa}
Un cero en la tabla casi siempre delata dependencia. Tomo $\fxy(10,7)=0$.
\paso{Comparo con el producto de sus m\'argenes}
\[f_X(10)\cdot f_Y(7)=0.30\,(0.10)=0.03\ne0=\fxy(10,7)\]
\paso{Concluyo}
Fall\'o en una celda, y con una basta. \res{\text{No son independientes}}
\end{pasos}
\end{resuelto}

\begin{turno}
Con la tabla del restaurante obt\'en la funci\'on condicional de $X$ dado $Y=9$ y despu\'es $P(X\ge9\mid Y=9)$.
\comprueba{se conoce $Y=9$: uso la columna $y=9$, que vale $0.05,\ 0.10,\ 0.20$ y suma $0.35$.
$f_{X\mid Y=9}(7)=\tfrac{0.05}{0.35}=\tfrac17$,\ \ $f_{X\mid Y=9}(9)=\tfrac{0.10}{0.35}=\tfrac27$,\ \ $f_{X\mid Y=9}(10)=\tfrac{0.20}{0.35}=\tfrac47$.
Entonces $P(X\ge9\mid Y=9)=\tfrac27+\tfrac47=\tfrac67\approx0.857$.}
\end{turno}

\subsection{Funci\'on de distribuci\'on conjunta (la tabla acumulada)}

\begin{idea}{$F_{XY}$ acumula hacia abajo y hacia la izquierda}
\[F_{XY}(x,y)=P(X\le x,\ Y\le y)\]
Cada celda de $F$ es la suma de todas las celdas de $f$ que est\'an \textbf{arriba y a la izquierda} de ella, incluy\'endola.
\end{idea}

\begin{resuelto}{Distribuci\'on conjunta de la f\'abrica}{Apuntes \textbullet{} Ej.~3.7 (en blanco en tus apuntes)}
\begin{pasos}
\paso{Primer rengl\'on ($y=0$): acumulo s\'olo hacia la derecha}
$F(0,0)=0.10$;\quad $F(1,0)=0.10+0.20=0.30$;\quad $F(2,0)=0.10+0.20+0.20=0.50$.
\paso{Segundo rengl\'on ($y=1$): acumulo los renglones $y=0$ y $y=1$}
$F(0,1)=0.10+0.04=0.14$;\quad $F(1,1)=0.10+0.20+0.04+0.08=0.42$;\quad
$F(2,1)=0.50+0.20=0.70$.
\paso{Tercer rengl\'on ($y=2$): acumulo los tres renglones}
$F(0,2)=0.10+0.04+0.06=0.20$;\quad $F(1,2)=0.20+0.40=0.60$;\quad $F(2,2)=1.00$.
\paso{Resultado}
\begin{center}
\begin{tabular}{|C|C|C|C|}\hline
\hd{$F_{XY}$} & \hx{$x=0$} & \hx{$x=1$} & \hx{$x=2$}\\\hline
\hy{$y=0$} & 0.10 & 0.30 & 0.50\\\hline
\hy{$y=1$} & 0.14 & \ce{0.42} & 0.70\\\hline
\hy{$y=2$} & 0.20 & 0.60 & 1.00\\\hline
\end{tabular}
\end{center}
\end{pasos}
\verifica{la \'ultima celda vale 1; el \'ultimo rengl\'on $(0.20,\,0.60,\,1.00)$ es la acumulada de $X$ sola y la \'ultima columna
$(0.50,\,0.70,\,1.00)$ es la acumulada de $Y$ sola.}
\enpalabras{la celda marcada dice que $P(X\le1,\,Y\le1)=0.42$.}
\end{resuelto}

\subsection{Construir la tabla desde el experimento}

A veces no te dan la tabla: hay que armarla. La receta es listar todos los resultados posibles, anotar cu\'anto valen
$X$ y $Y$ en cada uno y calcular su probabilidad.

\begin{resuelto}{Transistores defectuosos}{Serie 3 \textbullet{} problema 10 a}
\begin{enunciado}
Un dispositivo tiene cinco transistores, tres de ellos defectuosos. Se prueban uno por uno.
$X$ = n\'umero de pruebas hasta encontrar el primer defectuoso; $Y$ = n\'umero de pruebas \textbf{adicionales} hasta encontrar el segundo.
Determinar la funci\'on de probabilidad conjunta.
\end{enunciado}
\begin{pasos}
\paso{Entiendo las variables con un caso}
Llamo D a defectuoso y B a bueno. Si las pruebas salen B, D, B, D: el primer D apareci\'o en la prueba 2 ($X=2$)
y el segundo D apareci\'o 2 pruebas despu\'es ($Y=2$).
\paso{Listo todas las secuencias posibles hasta el segundo D}
S\'olo hay 2 buenos, as\'i que antes del segundo D puede haber 0, 1 o 2 buenos repartidos de seis maneras.
\paso{Calculo cada probabilidad multiplicando, sin reemplazo}
Cada factor es (los que quedan del tipo que sale) entre (los que quedan en total).
\begin{center}\small
\begin{tabular}{|>{\centering\arraybackslash}p{2.1cm}|C|C|p{5.6cm}|C|}\hline
\hd{Secuencia} & \hx{$X$} & \hy{$Y$} & \hd{C\'alculo} & \hd{Prob.}\\\hline
D D & 1 & 1 & $\tfrac35\cdot\tfrac24$ & $\tfrac{3}{10}$\\\hline
D B D & 1 & 2 & $\tfrac35\cdot\tfrac24\cdot\tfrac23$ & $\tfrac{2}{10}$\\\hline
D B B D & 1 & 3 & $\tfrac35\cdot\tfrac24\cdot\tfrac13\cdot1$ & $\tfrac{1}{10}$\\\hline
B D D & 2 & 1 & $\tfrac25\cdot\tfrac34\cdot\tfrac23$ & $\tfrac{2}{10}$\\\hline
B D B D & 2 & 2 & $\tfrac25\cdot\tfrac34\cdot\tfrac13\cdot1$ & $\tfrac{1}{10}$\\\hline
B B D D & 3 & 1 & $\tfrac25\cdot\tfrac14\cdot1\cdot1$ & $\tfrac{1}{10}$\\\hline
\end{tabular}
\end{center}
\porque{Ejemplo con D B D: al inicio hay 3 defectuosos de 5 ($\tfrac35$). Sacado uno, quedan 4 piezas y 2 son buenas ($\tfrac24$).
Sacado el bueno, quedan 3 piezas y 2 son defectuosas ($\tfrac23$).}
\paso{Acomodo en la tabla}
\begin{center}
\begin{tabular}{|C|C|C|C|C|}\hline
\hd{$x\,\backslash\,y$} & \hy{1} & \hy{2} & \hy{3} & \hx{$f_X(x)$}\\\hline
\hx{1} & 0.3 & 0.2 & 0.1 & \mx{0.6}\\\hline
\hx{2} & 0.2 & 0.1 & 0 & \mx{0.3}\\\hline
\hx{3} & 0.1 & 0 & 0 & \mx{0.1}\\\hline
\hy{$f_Y(y)$} & \my{0.6} & \my{0.3} & \my{0.1} & \textbf{1}\\\hline
\end{tabular}
\end{center}
\end{pasos}
\verifica{las seis probabilidades suman $\tfrac{10}{10}=1$: no falt\'o ninguna secuencia.
Esta tabla se vuelve a usar en el cap\'itulo \ref{cap:cov} para calcular la correlaci\'on.}
\end{resuelto}

\begin{turno}
(Ej.~3.1 de tus apuntes.) Se lanzan dos dados. $X_1$ = resultado del primero, $X_2$ = suma de los dos.
\textbf{a)} $\fxy(2,5)$\quad \textbf{b)} $\fxy(3,2)$\quad \textbf{c)} $f_{X_2}(4)$.
\comprueba{\textbf{a)} primer dado 2 y suma 5 obliga a que el segundo sea 3: un caso de 36, $\tfrac1{36}$.\quad
\textbf{b)} imposible (con un 3 la suma m\'inima es 4): $0$.\quad
\textbf{c)} suma 4 con $(1,3),(2,2),(3,1)$: $\tfrac{3}{36}$.}
\end{turno}
